%%=====================================================================
%%  preamble-oup.tex — the OUP port.
%%
%%  The class already provides: amsmath amssymb amsthm array caption
%%  graphicx hyperref multirow natbib rotating tikz xcolor url listings
%%  algorithm algpseudocode footmisc wrapfig tcolorbox subfloat.
%%  Do NOT reload any of them; several take options the class has set.
%%
%%  Dropped from the article-class preamble, because the class owns them:
%%    geometry fontenc inputenc lmodern microtype setspace  (page + fonts)
%%    natbib hyperref caption                               (refs + captions)
%%    the float-placement counters and \captionsetup        (class house style)
%%    the theorem environments                              (unused in the body)
%%  Kept from it: booktabs, threeparttable, mathtools, bm, float, and the
%%  whole macro block, which the sections depend on.
%%=====================================================================

% --- what the class does not supply ---------------------------------
\usepackage{booktabs}      % \cmidrule (16 uses) and \addlinespace (12); also
                           % supplies \bottomrule, which the class calls \botrule
\usepackage{threeparttable}% table notes under tables (30 uses)
\usepackage{mathtools}
\usepackage{bm}
\usepackage{float}
\usetikzlibrary{positioning,calc,fit,backgrounds,arrows.meta,shapes.geometric,decorations.pathreplacing,patterns}
\graphicspath{{figures/}}
\newcommand{\societylogo}{}  % the template defines this in its preamble

% --- cross-reference helpers (cleveref is not installed here) --------
\newcommand{\secref}[1]{Section~\ref{#1}}
\newcommand{\secsref}[2]{Sections~\ref{#1} and~\ref{#2}}
\newcommand{\tabref}[1]{Table~\ref{#1}}
\newcommand{\figref}[1]{Figure~\ref{#1}}
\newcommand{\eqnref}[1]{Equation~\eqref{#1}}
\newcommand{\appref}[1]{Appendix~\ref{#1}}
\newcommand{\panelref}[2]{Panel~#1 of \tabref{#2}}

% --- drafting helpers (delete before submission) ---------------------
\newcommand{\TODO}[1]{\textcolor{red}{[\textbf{TODO:} #1]}}
\newcommand{\PH}[1]{\textcolor{gray}{\ensuremath{\langle}\text{#1}\ensuremath{\rangle}}}
\newcommand{\PENDGROSS}{\TODO{pending gross-alignment verification}}
\newcommand{\PROVEN}{\textsc{[proven]}}
\newcommand{\SKETCH}{\textsc{[sketch]}}
\newcommand{\CONJ}{\textsc{[conjecture]}}

%%---------------------------------------------------------------------
%%  NOTATION — identical to the article-class build; the sections use it.
%%---------------------------------------------------------------------
\newcommand{\N}{N}
\newcommand{\Tw}{T}
\newcommand{\K}{K}
\newcommand{\Rmat}{\mathbf{R}}
\newcommand{\Cmat}{\mathbf{C}}
\newcommand{\Wmat}{\mathbf{W}}
\newcommand{\Dmat}{\mathbf{D}}
\newcommand{\Proj}{\mathbf{P}}
\newcommand{\Imat}{\mathbf{I}}
\newcommand{\rt}{\tilde{\mathbf{r}}}
\newcommand{\Zmat}{\mathbf{Z}}
\newcommand{\Lam}{\bm{\Lambda}}
\newcommand{\ones}[1]{\bm{1}_{#1}}
\newcommand{\Pgics}{\mathbf{P}_{\mathrm{GICS}}}
\newcommand{\SR}{\mathrm{SR}}
\newcommand{\Exp}{\mathbb{E}}
\newcommand{\Prob}{\mathbb{P}}
\newcommand{\Var}{\mathrm{Var}}
\newcommand{\Cov}{\mathrm{Cov}}
\newcommand{\real}{\mathbb{R}}
\newcommand{\tr}{\operatorname{tr}}
\newcommand{\rank}{\operatorname{rank}}
\DeclareMathOperator*{\argmin}{arg\,min}
\DeclareMathOperator*{\argmax}{arg\,max}
\newcommand{\Part}{\mathcal{P}}
\newcommand{\PartOne}{\Part_{\mathrm{one}}}
\newcommand{\PartRnd}{\Part_{\mathrm{rand}}}
\newcommand{\alphaFF}{\alpha_{\mathrm{FF5+MOM}}}
\newcommand{\ARI}{\mathrm{ARI}}
\newcommand{\nogroup}{\emph{no grouping}}
\newcommand{\randgroup}{\emph{random grouping}}
